\documentclass[10pt,a4paper]{amsart}
\usepackage[margin=19mm]{geometry}
\usepackage{amsmath,amssymb,mathtools}
\usepackage[scr=rsfs,cal=euler]{mathalfa}
\usepackage{xfrac}
\usepackage[unicode,hidelinks,bookmarks=false]{hyperref}
\newcommand{\ie}{i.e.\ }
\newcommand{\cf}{cf.\ }
\newcommand{\resp}{respectively\ }
\newcommand{\Subsection}{\subsection}
\providecommand{\nobreakdash}{\nobreak\hbox{-}}
\setlength{\emergencystretch}{2em}
\multlinegap0pt
\usepackage[shortlabels]{enumitem}
\usepackage{mathtools}
\usepackage{amssymb}
\usepackage{float}
\usepackage{caption}
\usepackage{tikz}
\usepackage{algorithm}\usepackage{algpseudocode}
\usepackage{xcolor}
\usepackage{multirow}
\usepackage{bm}
\newtheorem{lemma}{Lemma}
\usepackage{cleveref}
\crefname{lemma}{Lemma}{Lemmas}
\newcommand{\K}{\mc{K}}
\newcommand{\Z}{\mathbb{Z}}
\newcommand{\abs}[1]{|{#1}|}
\newcommand{\mc}{\mathcal}
\newcommand{\torsion}[1]{\bar{#1}^{\,*}}
\title{{\textbf{\large{Generators of $H^1(\Gamma, \partial \Gamma^c)$ with $\partial \Gamma^c \subset \partial \Gamma$ for Triangulated Surfaces $\Gamma$: Construction and Classification of Global Loops}}}}
\author{Silvano Pitassi}
\date{Source: arXiv:2504.05124v2 (CC BY 4.0)}
\begin{document}
\maketitle

\noindent\emph{Source and licence.} {\textbf{\large{Generators of $H^1(\Gamma, \partial \Gamma^c)$ with $\partial \Gamma^c \subset \partial \Gamma$ for Triangulated Surfaces $\Gamma$: Construction and Classification of Global Loops}}}. Version arXiv:2504.05124v2 (CC BY 4.0), distributed by arXiv under the Creative Commons Attribution 4.0 International licence, http://creativecommons.org/licenses/by/4.0/, as recorded in the arXiv licence metadata for this exact version and shown on the abstract page https://arxiv.org/abs/2504.05124v2. Full licence notice: \texttt{licenses/arxiv-2504.05124-CC-BY-4.0.txt}.\par\smallskip

\noindent\textit{Editorial scope.} Counted unit: the article's lemma lem:homology (source0.tex lines 743--759). The article's complete original proof of it is delivered with it (source0.tex lines 760--780, one \texttt{proof} environment of 21 lines). No mathematics is written, repaired or re-derived: the statement and the proof are the article's own lines, in the article's own notation, and no retained line is substituted or re-flowed.  Retained uncounted as the source-local context the proof needs: lines 726--729.  Nothing was omitted from any retained span, so the package carries no omission record.  No retained line contains a citation, so the wrapper carries no bibliography and no bibliography entry.  No retained line contains a figure environment, an image-inclusion command or drawing code, so no image is delivered and the package declares no asset dependency.  Editorial formatting only, in addition: an AMS \texttt{amsart} preamble that restates the article's own declarations and macros used by the retained lines, namely the environments \texttt{lemma} on separate counters, together with the standard packages those lines need.  The retained environments print in this document as Lemma 1 in order of appearance, whereas the article prints the same items with the numbering of its own full document.  The complete native source of the article as distributed in the arXiv e-print for this exact version is bundled unchanged as \texttt{sources/176117/source0.tex}.
\begin{equation}
\partial^M_2 f = \sum_{e \in E_M} \iota^M(f, e) \, e,
\label{eq:expansion}
\end{equation}

\begin{lemma}[Homology characterization of $\bar \K$]
\label{lem:homology}
For each $e \in E_M$, let $c_e$ be the unique cycle in the graph $T \cup \{e\}$, naturally oriented according to the orientation induced by $e$.
Note that $c_e$ can considered as a $1$-cycle in $\bar{\K}$.
Let $E_M = E_M^{\rm I} \sqcup E_M^{\rm II}$, and denote by $g := \abs{E_M} = 2 - \chi$ the genus of $\bar \K_M$.
We distinguish between the following two cases:
\begin{enumerate}[label=(\roman*)]
\item If $\abs{E_M^{\rm II}} = 0$, then $H_1(\bar \K; \Z) \cong \Z^g$.
\item If $\abs{E_M^{\rm II}} > 0$, then $H_1(\bar \K;\Z)\cong\Z^{g-1} \oplus \Z/2\Z$.
In this case, the torsion generator $\torsion c$ is given by 
\begin{equation}
\torsion c := \sum_{e \in E_M^{\rm{II}}} \eta_e\, c_e,
\label{eq:torsion}
\end{equation}
where each $\eta_e \in \{-1,+1\}$ is the sign of the coefficient $\iota^M(f,e)$ in the expression of $\partial^M_2 f$ as a linear combination of $1$-cells $e \in E_M^{\rm{II}}$ in \cref{eq:expansion}.
\end{enumerate}
\end{lemma}

\begin{proof}
Since $\bar{\K}_M$ has only one $0$-cell, every $1$-cell $e \in E_M$ forms a $1$-cycle.

In the first case, when $\abs{E_M^{\mathrm{II}}} = 0$, we have $E_M = E_M^{\mathrm{I}}$.
Therefore, each $1$-cell $e \in E_M$ is a homology generator of $H_1(\bar{\K}_M; \mathbb{Z})$, since $\partial_2^M f \neq e$.
Using the fact that $\bar{\K}$ is homotopy equivalent to $\bar{\K}_M$, we conclude that each $1$-cycle $c_e$ is also a homology generator of $H_1(\bar{\K}; \mathbb{Z})$.

In the second case, where $\abs{E_M^{\mathrm{II}}} > 0$, fix a $1$-cell $e^* \in E_M^{\mathrm{II}}$.
The boundary operator $\partial_2^M f$ satisfies
\[
\partial_2^M f = \sum_{e \in E_M^{\mathrm{II}}} 2 \eta_e e \eqqcolon 2c^*,
\]
where $c^*$ is the (unique) torsion generator of $\bar{\K}_M$.

For each $e \in E_M \setminus \{e^*\}$, the $1$-cell $e$ is a homology generator of $H_1(\bar{\K}_M; \Z)$, since $\partial_2^M f \neq e$.
This means there are $g - 1$ independent homology generators.

Next, let $c_e$ be the unique $1$-cycle of $\bar{\K}$ associated with each $e \in E_M$, as in the case where $\abs{E_M^{\mathrm{II}}} = 0$, and define the $1$-cycle $\torsion c$ as in \cref{eq:torsion}.
Using again the fact that $\bar{\K}$ is homotopy equivalent to $\bar{\K}_M$, each $c_e$ for $e \in E_M \setminus \{e^*\}$ is a homology generator of $H_1(\bar{\K}; \mathbb{Z})$.
Furthermore, $\torsion c$ is the torsion generator of $H_1(\bar{\K}; \mathbb{Z})$, since $\torsion c$ is homotopy equivalent to the torsion $1$-cycle $c^*$ in $\bar{\K}_M$.
\end{proof}

\end{document}
